\hypertarget{appendix-37856d997bb0}{}
\section*{Reference binding and converted temporaries}
\section{Reference binding and converted temporaries}\label{reference-binding-and-converted-temporaries}

\textbf{Curriculum depth:} Introduced · \textbf{First seen:} C008 · \textbf{Latest development:} C008

\subsection{Definition and mechanism}\label{definition-and-mechanism}

Binding chooses an initial referent under type/category rules. A const int\& bound to an int lvalue tracks that integer; a const double\& initialized from the same int binds a converted temporary snapshot. Base-subobject binding need not copy; bit-fields cannot provide ordinary mutable integer references.

\subsection{Prerequisites}\label{prerequisites}

\begin{itemize}
\tightlist
\item
  \hyperref[appendix-d6d562ca9fbb]{References and aliases} --- The initial alias and its referent determine later binding and deduction behavior.
\item
  \hyperref[appendix-1121b30eaeb0]{Value categories and materialization} --- Expression category determines viable binding and whether temporary identity is needed.
\end{itemize}

\subsection{Used by}\label{used-by}

\begin{itemize}
\tightlist
\item
  \hyperref[appendix-a978b29bbbb0]{Temporary lifetime extension}
\end{itemize}

\subsection{Recall and next depth}\label{recall-and-next-depth}

Mutate the original after both bindings and predict their values. Pending: user-defined conversion overloads and volatile binding. Written coverage does not establish learner mastery.

\subsection{Source and navigation}\label{source-and-navigation}

\href{Continuous\%20Notes/Computing\%20Foundations\%20-\%20From\%20Bits\%20to\%20C++\%20and\%20Systems\%20APIs\#Binding\%20chooses\%20a\%20referent,\%20not\%20always\%20the\%20one\%20you\%20imagined}{Teaching explanation and examples} · \hyperref[appendix-fd090d982a15]{Exact coverage} · \href{}{Dependency Map (02-Foundations-II, appendix-4ad47bdd269a)}
